\section*{الفصل \ref{ch:b1:proteins} --- الحموض الأمينية والبروتينات}

\begin{solution}{exo:b1:proteins:1}
$\mathrm{^+H_3N{-}CH(R){-}COO^-}$: أي \omterm{def:b1:proteins:aminoacid}{شاردة ثنائية القطب}. وLeu غير
\omterm{def:b1:water-small-molecules:water}{قطبي} أليفاتي؛ وSer \omterm{def:b1:water-small-molecules:water}{قطبي} غير مشحون؛ وLys موجب الشحنة؛ وGlu سالب
الشحنة؛ وPhe عطري.
\end{solution}

\begin{solution}{exo:b1:proteins:2}
الأولية: التسلسل، وروابط ببتيدية. والثانوية: طي موضعي للهيكل (لولب،
وصفيحة)، \omterm{def:b1:water-small-molecules:water}{وروابط هيدروجينية} في الهيكل. والثالثية: طية السلسلة كلها؛
\omterm{prop:b1:water-small-molecules:properties}{والأثر الكاره للماء}، \omterm{def:b1:water-small-molecules:water}{والروابط الهيدروجينية}، والجسور الملحية، وتماسات
\omterm{def:b1:water-small-molecules:bonds}{فان دير فالس}، \omterm{def:b1:proteins:tertiary}{والجسور ثنائية الكبريت}. والرابعية: تركيب الوحدات،
بالقوى غير التساهمية نفسها.
\end{solution}

\begin{solution}{exo:b1:proteins:3}
$450\times 110 = \qty{50}{kDa}$. واللولب: $450\times 0.15 = \qty{68}{nm}$.
والخصلة: $450\times 0.35 = \qty{158}{nm}$.
\end{solution}

\begin{solution}{exo:b1:proteins:4}
نحو $0.98$ و$0.76$: أي يُسلَّم خُمس الحمولة (\qty{22}{\%}).
\end{solution}

\begin{solution}{exo:b1:proteins:5}
الزمر: الأمين الطرفي (9.5)، \omterm{def:b1:proteins:aminoacid}{والسلسلة الجانبية} عند Lys (10.5)،
\omterm{def:b1:proteins:aminoacid}{والسلسلة الجانبية} عند Asp (3.9)، والكربوكسيل الطرفي (2). وعند pH 7:
$+1 + 1 - 1 - 1 = 0$. وعند pH 2: $+1 + 1 + 0 - 0.5 = +1.5$
(فالكربوكسيل الطرفي مشرَّد بنصفه). وعند pH 12:
$0 + 0 - 1 - 1 = -2$ (\omterm{def:b1:water-small-molecules:ph}{فحمض} Lys منزوع البرتنة في معظمه). والشحنة صفر بين
p$K_a$ عند Asp (3.9) وعند الزمرة الأمينية (9.5)، حيث تكون الشحنة
$0$ على طول المدى: فتكون p$I \approx 6.7$ (متوسط 3.9 و9.5).
\end{solution}

\begin{solution}{exo:b1:proteins:6}
\omterm{def:b1:proteins:aminoacid}{السلسلة الجانبية} للبرولين مرتبطة راجعةً إلى نتروجينه: فلا يحمل
النتروجين هيدروجينًا يمنحه \omterm{def:b1:water-small-molecules:water}{للرابطة الهيدروجينية} في اللولب، و$\phi$
محبوسة عند نحو $-60^\circ$، وهذا لا يُحتمل إلا عند الدورة الأولى من
لولب. أما الغليسين، وسلسلته الجانبية هيدروجين، فلا يعاني أي تصادم
فراغي ويستطيع أن يتخذ قيم $\phi$ و$\psi$ ممنوعة على كل بقية أخرى، بما
فيها قيم الانعطافات الضيقة.
\end{solution}

\begin{solution}{exo:b1:proteins:7}
اقترانات ثمانية سيستئينات: $7\times 5\times 3\times 1 = 105$؛ فإعادة
الأكسدة العشوائية تعطي المجموعة الأصلية نحو \qty{1}{\%} من المرات.
واسترداد النشاط الكامل تقريبًا بيّن أن السلسلة، لا المصادفة، هي التي
اختارت الاقتران: فالتسلسل يشفّر الطية \omterm{def:b1:proteins:tertiary}{والجسور ثنائية الكبريت} لا تفعل
إلا أن تقفلها.
\end{solution}

\begin{solution}{exo:b1:proteins:8}
$P_{50}^{2.8} = 33.4$. وعند 2: $2^{2.8} = 6.96$، فيكون $Y = 0.17$؛
وعند 3.5: $0.50$؛ وعند 8: $8^{2.8} = 337$، فيكون $Y = 0.91$. ومع
$n = 1$: $Y =
P/(3.5 + P)$: أي $0.36$ و$0.50$ و$0.70$. والتسليم بين
8 و2: $0.74$ عند التعاوني، و$0.34$ عند غير التعاوني: \omterm{def:b1:proteins:peptide}{فالبروتين}
التعاوني يسلّم ضعفًا.
\end{solution}

\begin{solution}{exo:b1:proteins:9}
رباعي من أربع وحدات متطابقة كتلة كل منها \qty{50}{kDa}. وأجرِ هلام
SDS بلا عامل مرجع: فإن ظهرت حزم عند 100 أو 150 أو \qty{200}{kDa}
\omterm{def:b1:proteins:tertiary}{فالجسور ثنائية الكبريت} تربط الوحدات؛ وإن بقيت حزمة \qty{50}{kDa}
وحدها فهي مرتبطة ارتباطًا غير تساهمي.
\end{solution}

\begin{solution}{exo:b1:proteins:10}
الغلوتامات مشحونة مماهة عند السطح؛ أما الفالين فكاره للماء ويحدث بقعة
لاصقة يستبعدها الماء. وفي الحالة T منزوعة الأكسجين ينكشف جيب كاره
للماء متكامل معها على جزيء مجاور، فتلتصق الجزيئات طرفًا بطرف في ألياف
تشوه الكرية الحمراء. وفي الرئتين تخفي الحالة R الجيب فتذوب الألياف؛
وفي \omterm{def:b1:body-plans-tissues:tissue}{الأنسجة}، حيث يكون الهيموغلوبين منزوع الأكسجين، تتكون --- فتتمنجل
\omterm{def:b1:cell-unit-of-life:cell}{الخلايا} وتسد الأوعية الصغيرة هناك.
\end{solution}

\begin{solution}{exo:b1:proteins:11}
عند \qty{37}{\celsius} يكون $RT = \qty{2.58}{kJ/mol}$: فيكون
$K = e^{-40/2.58} =
e^{-15.5} = \num{1.8e-7}$: أي جزيء واحد من خمسة
ملايين مفكوك. وعند \qty{45}{\celsius} مع \qty{10}{kJ/mol} يكون
$RT = 2.64$: فيكون $K = e^{-3.8} =
0.022$: أي \qty{2}{\%} مفكوكة،
وفي ارتفاع حاد. والاستقرار الحدي يتيح \omterm{def:b1:proteins:peptide}{للبروتينات} أن تغير شكلها وهي
تعمل (\omterm{def:b1:proteins:allostery}{الألوستيرية}، والحفز)، وأن تُفك لنقلها عبر \omterm{def:b1:membranes-transport:membrane}{الأغشية}، وأن تُهدم
وتُستبدل سريعًا حين تتضرر أو تنتفي الحاجة إليها.
\end{solution}

\begin{solution}{exo:b1:proteins:12}
بيّن أنفنسن أن التسلسل يكفي لتحديد الطية، وتبين \omterm{def:b1:proteins:allostery}{الألوستيرية} أن الشكل
--- وتغيراته --- هو الوظيفة. لكن \omterm{def:b1:cell-unit-of-life:cell}{الخلية} تتدخل في كل خطوة: \omterm{def:b1:proteins:peptide}{فالبروتينات}
المرافقة تنقذ سلاسل كانت لتتكتل قبل الانطواء، لأن العصارة المزدحمة
ليست أنبوب اختبار مخففًا؛ واستبدال واحد (المنجلية) يعطي بروتينًا مطويًا
طيًا صحيحًا لكن سطحه خاطئ، فلا بد أن يشمل «الشكل» كيمياء السطح لا
الهيكل وحده؛ وكثير من \omterm{def:b1:proteins:peptide}{البروتينات} يحتاج إلى تعديلات أو شركاء أو ربائط
تُضاف بعد الطي قبل أن يعمل. فالجملة تذكر المبدأ؛ \omterm{def:b1:cell-unit-of-life:cell}{والخلية} توفر الشروط.
\end{solution}

\begin{solution}{pb:b1:proteins:1}
\textbf{1.} $150\times 1.34 = \qty{201}{mL}$ من $\mathrm{O_2}$ في اللتر.
\textbf{2.} $e^{3.51} = 33.4$؛ و$e^{4.67} = 106.7$؛
و$e^{7.25} = 1408$.
\textbf{3.} الرئتان: $1408/(33.4 + 1408) = 0.977$. \omterm{def:b1:body-plans-tissues:tissue}{والأنسجة}:
$106.7/(33.4 + 106.7) = 0.762$.
\textbf{4.} $0.977 - 0.762 = 0.215$: أي $0.215\times 201 = \qty{43}{mL}$
في اللتر.
\textbf{5.} $250/43 = \qty{5.8}{L/min}$ --- وهو قريب من الصبيب
القلبي في الراحة.
\textbf{6.} الميوغلوبين: في الرئتين $13.3/13.67 = 0.973$؛ وفي
\omterm{def:b1:body-plans-tissues:tissue}{الأنسجة} $5.3/5.67 = 0.935$؛ فيسلّم $0.038\times 201 = \qty{7.7}{mL}$
في اللتر، أي أقل بست مرات: فالحامل الزائدي الذي يحمّل جيدًا لا يستطيع
التفريغ عند ضغوط \omterm{def:b1:body-plans-tissues:tissue}{الأنسجة}.
\textbf{7.} \omterm{def:b1:proteins:allostery}{التعاونية} تضع الجزء الحاد من المنحني بين ضغطي العمل،
فيطلق هبوط صغير في الضغط كسرًا كبيرًا من الحمولة.
\textbf{8.} $e^{2.78} = 16.1$؛ فيكون $Y = 16.1/(33.4 + 16.1) = 0.325$.
\textbf{9.} $e^{4.21} = 67.4$؛ فيكون $Y = 16.1/(67.4 + 16.1) = 0.193$.
\textbf{10.} بلا أثر بور: $(0.977 - 0.325)\times 201 = \qty{131}{mL}$؛
ومعه: $(0.977 - 0.193)\times 201 = \qty{158}{mL}$ في اللتر؛ والكسب
\qty{27}{mL}، أي \qty{20}{\%}.
\textbf{11.} $2.7/3.07 = 0.88$ مشبع؛ وعند \qty{0.5}{kPa} يكون
$0.5/0.87 = 0.57$: فيطلق ثلث مخزونه إلى \omterm{def:b1:eukaryotic-cell:mitochondrion}{المتقدرات} في أثناء التقلص،
ويعيد التحميل بين التقلصات.
\textbf{12.} ترتبط البروتونات في مواقع (الهيستيدينات، والنهايات
الأمينية) بعيدة عن الهيمات، وترسّخ بارتباطها الحالة T، فتخفض الألفة
عند الهيمات: أي ربيطة في موقع تغيّر الألفة في موقع آخر --- وهي
\omterm{def:b1:proteins:allostery}{الألوستيرية}.
\textbf{13.} $e^{5.45} = 233$؛ فالرئتان $233/(33.4 + 233) = 0.875$؛
\omterm{def:b1:body-plans-tissues:tissue}{والأنسجة} $0.762$؛ والتسليم $0.113\times 201 = \qty{23}{mL}$: أي نحو
نصف ما عند مستوى سطح البحر.
\textbf{14.} $e^{3.88} = 48.4$: فالرئتان $233/(48.4 + 233) = 0.828$؛
\omterm{def:b1:body-plans-tissues:tissue}{والأنسجة} $106.7/(48.4 + 106.7) = 0.688$؛ والتسليم
$0.140\times 201 =
\qty{28}{mL}$.
\textbf{15.} يهبط التحميل قليلًا (من 0.875 إلى 0.828) لكن التفريغ
يهبط أكثر (من 0.762 إلى 0.688)، فينمو الفرق: إذ ما تزال الرئتان عند
\qty{7}{kPa} على القمة المستوية للمنحني بينما \omterm{def:b1:body-plans-tissues:tissue}{الأنسجة} على جزئه الحاد.
\textbf{16.} السعة $180\times 1.34 = \qty{241}{mL}$؛
و$0.140\times
241 = \qty{34}{mL}$ في اللتر --- أي استعادة ثلاثة أرباع
ما عند مستوى سطح البحر.
\textbf{17.} $6^{2.8} = e^{2.8\times 1.79} = e^{5.02} = 151$:
فالرئتان $233/384 = 0.607$، \omterm{def:b1:body-plans-tissues:tissue}{والأنسجة} $106.7/258 = 0.414$: والتسليم
$0.193$، وهو أفضل على الورق --- لكن الدم الشرياني لن يكون مشبعًا إلا
\qty{61}{\%}، وأي هبوط آخر في ضغط الرئة أو أي طلب لضغط نسيجي أعلى
كان ليترك الدماغ في عوز؛ ولذلك يقف التعديل الحقيقي قرب
\qty{4}{kPa}.
\textbf{18.} $48.4/(33.4 + 48.4) = 0.59$.
\textbf{19.} $e^{2.68} = 14.6$؛ و$48.4/(14.6 + 48.4) = 0.77$.
\textbf{20.} عند الضغط نفسه يرتبط \omterm{def:b1:proteins:peptide}{بروتين} الجنين أكثر: فبينما يفرغ دم
الأم نحو \qty{59}{\%} ويحمّل دم الجنين نحو \qty{77}{\%}، يجري
الأكسجين هابطًا مع تدرج الضغط الذي يبقيه فرق الألفتين عبر الغشاء
المشيمي.
\textbf{21.} يرتبط BPG بالحالة T ويخفض الألفة؛ وسلاسل $\gamma$
ترتبط به ارتباطًا ضعيفًا، فيبقى هيموغلوبين الجنين أقرب إلى ألفته
الذاتية العالية: أي $P_{50}$ أدنى.
\textbf{22.} الرئتان: $1408/(14.6 + 1408) = 0.990$؛ \omterm{def:b1:body-plans-tissues:tissue}{والأنسجة}
$106.7/(14.6
+ 106.7) = 0.880$؛ والتسليم
$0.110\times 201 = \qty{22}{mL}$، أي نصف تسليم البالغ: وهو جيد لأخذ
الأكسجين من الأم، رديء لتسليمه بعد أن تعمل الرئتان، ومن هنا الانتقال.
\textbf{23.} $5^{2.8} = e^{2.8\times 1.609} = e^{4.51} = 90.6$:
فالرئتان $1408/1499 = 0.94$، \omterm{def:b1:body-plans-tissues:tissue}{والأنسجة} $106.7/197 = 0.54$: والتسليم
$0.40$، أي نحو ضعف الطبيعي --- والشخص مزرق الشفتين قليلًا لكنه يحتمل
ذلك جيدًا؛ والثمن احتياطي أصغر على الارتفاع.
\textbf{24.} $Y = P/(3.5 + P)$: فالرئتان $0.79$، \omterm{def:b1:body-plans-tissues:tissue}{والأنسجة} $0.60$:
والتسليم $0.19\times 201 = \qty{38}{mL}$ --- وفي العضلة التي تعدو عند
\qty{2.7}{kPa} يكون $0.44$: والتسليم $0.35\times 201 = \qty{70}{mL}$
مقابل \qty{158}{mL}. فبلا \omterm{def:b1:proteins:allostery}{تعاونية} يحمّل الحامل تحميلًا رديئًا ويفرغ
تفريغًا رديئًا.
\textbf{25.} \qty{43}{mL} من $\mathrm{O_2}$ في اللتر في الراحة،
و\qty{158}{mL} في اللتر في عضلة تعدو؛ والفرق يأتي من \omterm{def:b1:proteins:allostery}{التعاونية}
(المنحني السيني) ومن أثر بور (إزاحة $P_{50}$ في \omterm{def:b1:water-small-molecules:ph}{الحمض}).
\end{solution}
